% Code of Conduct, officer variant.
%
% NOT A WORD-SWAP OF THE INTERN CODE. Four sections are substantively DIFFERENT, and two of
% them are inverted -- copying the intern code and changing the noun would have produced a
% document that is wrong in the officer's favour in one place and against them in another:
%
%   4  Availability      an officer has no supervisor and no recorded weekly hours.
%   5  Outside work      an unpaid intern is free to work anywhere and Paridyum says so.
%                        An officer holds a position of trust; the same sentence would be
%                        a much larger concession than anyone intends.
%   7  Speaking publicly INVERTED. An intern may not speak for Paridyum. An officer can,
%                        and the useful rule is about care rather than prohibition.
%   9  Breach            an officer is not "terminated" -- an office is vacated or removed
%                        by the directors. There is no Officership Agreement to point at.
%
% Section 8 also changes: the Grievance Policy is an intern document (classes: [intern]),
% so it cannot be the route named here.
%
% WORDING IS A DRAFT FOR REVIEW BY QUALIFIED COUNSEL. That was true of the intern code and
% it is more true here, because the duties an officer owes are not ones this estate has
% written down anywhere yet.
% adapted-from: templates/policy/code-of-conduct.tex@6c9096ed039beec4de55d37fb3fdad40f4136a9c3c392bb30e5530bedd846f94
%
% NOT A SHARED FILE, AND NOT A COPY THAT DRIFTS SILENTLY. Of the sections in this pair,
% more than half differ in substance -- two of them inverted -- so there is no shared
% clause to extract. What there is instead is a risk: the intern document is amended and
% this one is not, and nothing notices.
%
% The digest above is of the file this was adapted from. `cargo test --test repository`
% recomputes it and fails when it moves, naming this file. That is the whole mechanism:
% it does not compare content -- the two are SUPPOSED to differ -- it makes it impossible
% to change one half without being told to review the other.
%
% When the intern document legitimately changes: review this one, make the matching
% change or decide it does not apply, then update the digest in the same commit.

\documentclass{paridyum}
\begin{document}

\paridyumletterhead
\paridyumdateref{\dkDocDate}{\dkDocRef}
\paridyumtitle{Code of Conduct (Officer)}

\textit{This code applies to everyone working at or with Paridyum. This copy is the one
given to an officer, and it differs from the copy given to an intern where holding office
makes a difference --- chiefly outside interests, speaking publicly, and what happens on a
breach.}

\section*{1. What This Covers}
This code sets out how Paridyum expects people to conduct themselves. It sits
alongside, and does not replace, the Information Security \& Acceptable Use
Policy, the Policy on Prevention of Sexual Harassment, and the Non-Disclosure
Agreement.

\section*{2. Treating People Decently}
Everyone at Paridyum is entitled to be treated with respect. The following are
not acceptable, whether in person, in writing, or in any online space used for
work:
\begin{itemize}[leftmargin=*, noitemsep]
    \item bullying, intimidation, or persistent unwelcome behaviour;
    \item discrimination or harassment on the basis of sex, gender, caste,
          religion, disability, sexual orientation, or any other personal
          characteristic;
    \item abusive, demeaning, or threatening language; and
    \item deliberately excluding a colleague from work they should be part of.
\end{itemize}
Sexual harassment is dealt with separately and more specifically by the Policy
on Prevention of Sexual Harassment, which sets out how to complain and to whom.

\section*{3. Honesty About Work}
The Officer agrees to:
\begin{itemize}[leftmargin=*, noitemsep]
    \item represent their own progress accurately, including when work is
          behind or has gone wrong;
    \item say when they do not know something, rather than guess in a way
          that others will rely on;
    \item attribute work honestly, including work produced with the
          assistance of AI tools; and
    \item raise a problem early rather than let it become someone else's
          surprise.
\end{itemize}
Reporting a mistake promptly is treated as doing the job properly, not as a
failure.

\section*{4. Availability and Communication}
An officer has no recorded weekly hours and no supervisor. What is expected instead is
that the other officers can reach them on a matter that needs an officer, and that they
say so in advance when they will be unreachable for any extended period --- particularly
where they are the signatory a document is routed to.

\section*{5. Outside Work and Conflicts of Interest}
\textbf{An officer may hold other work or interests, and holding office is not itself a
claim on their time.} No remuneration is paid for the office. But an officer is in a
position of trust, and the standard here is higher than the one that applies to an unpaid
intern: the question is not only whether outside work competes for time, but whether it
could compete with the judgement the office requires.

The Officer must tell the other officers if any of the following applies, in writing,
before rather than after:
\begin{itemize}[leftmargin=*, noitemsep]
    \item they work for, or are engaged by, a business that competes directly
          with a Paridyum product they are working on;
    \item they have a financial interest in a supplier, customer, or
          competitor of Paridyum;
    \item a close family member is in a position to be affected by a decision
          the Officer takes for Paridyum; or
    \item they are subject to an obligation to another party that would
          conflict with the Intellectual Property Assignment Agreement.
\end{itemize}
Disclosing a conflict is not a problem. Not disclosing one is. Where a conflict touches a
decision, the Officer takes no part in that decision --- and, where the decision is a
document, is not the signatory it routes to.

\section*{6. Company Property and Resources}
Paridyum accounts, systems, and equipment are provided for Paridyum work.
Incidental personal use is acceptable; using them for another business, or in a
way that puts Paridyum's information at risk, is not. The Information Security
\& Acceptable Use Policy sets the detail.

\section*{7. Speaking About Paridyum}
\textbf{An officer may speak on Paridyum's behalf, and this section is about care rather
than permission.} An officer's public statement is read as the company's whether or not
it was meant that way, so:
\begin{itemize}[leftmargin=*, noitemsep]
    \item claims about what Paridyum has built, operates or has been paid for must be
          true at the time they are made, and not merely intended;
    \item nothing confidential and no unreleased work is disclosed, whatever the
          setting; and
    \item a commitment that would bind Paridyum --- to a customer, a supplier or a
          candidate --- is made in writing, not in passing.
\end{itemize}
An officer's name and office are published on Paridyum's website, which is what makes
the first of those matter more than it would for anyone else.

\section*{8. Raising a Concern}
The Grievance Policy is an intern document and does not set the route here. A concern
about an officer's conduct is raised with any other officer, in writing, and is
considered by the officers other than the one it concerns --- who takes no part in it.
A concern raised in good faith will not be held against the person raising it, and that
holds whether it is raised by an officer or about one.

\section*{9. Breach}
A serious or repeated breach of this code may lead to removal of system access, and to
the Officer being asked to vacate the office or removed from it by the directors. An
office is not terminated in the way an engagement is; it is vacated or removed, and the
signatory register records the date either way. Obligations under the Non-Disclosure
Agreement and the Intellectual Property Assignment continue regardless.

\section*{10. Acknowledgement}
By signing below, the Officer confirms that they have read and understood this
Code of Conduct and agree to follow it.

\vspace{1.5em}
\paridyumsignblock

\vspace{1.5em}
\paridyumcountersign{Officer}{\dkPIdentityLegalName}

\end{document}
